\begin{table}[!htb]
\centering
\caption{Evaluasi Alternatif Distribusi Kubernetes \textit{Open-Source}}
\label{tab:k8s_distribution_comparison}
\footnotesize
\setlength{\tabcolsep}{3pt}
\begin{tabular*}{\textwidth}{@{\extracolsep{\fill}}|>{\raggedright\arraybackslash}m{2.9cm}|>{\centering\arraybackslash}m{1.5cm}|>{\centering\arraybackslash}m{1.5cm}|>{\centering\arraybackslash}m{1.5cm}|>{\centering\arraybackslash}m{1.5cm}|>{\centering\arraybackslash}m{2.5cm}|>{\centering\arraybackslash}m{0.9cm}|}
\hline
\multicolumn{1}{|>{\centering\arraybackslash}m{2.9cm}|}{\textbf{Distribusi}} & \textbf{\textit{Footprint} Ringan} & \textbf{Siap Produksi} & \textbf{Tanpa Dependensi} & \textbf{Eko\-sistem Add-on} & \textbf{Lisensi (Yayasan)} & \textbf{Total} \\
\hline
\textit{k3s} & 2 & 2 & 2 & 2 & Apache 2.0 (CNCF) & 8 \\
\hline
\textit{k0s} & 2 & 2 & 2 & 1 & Apache 2.0 (mandiri) & 7 \\
\hline
MicroK8s & 1 & 2 & 1 & 2 & Apache 2.0 (mandiri) & 6 \\
\hline
\textit{kind} & 1 & 0 & 1 & 0 & Apache 2.0 (CNCF) & 2 \\
\hline
\end{tabular*}
\skorlegendtools
\end{table}
